% tcy-star-error.tex -- 縦中横の方向契約のエラーゲート。
% 契約:
%   * yoko 文書で非 star 形を使うと ClassError（"tatechuyoko is not
%     allowed in yoko mode"）。
%   * yoko 文書で star 形は中身をそのまま通す（透過）。
%   * 未知 box の方向照会 \ptj@boxdir は ClassError（無告知の
%     yoko fallback はしない）。
% 本 fixture は yoko 文書で非 star 形を「エラーになることを確認」する:
% LaTeX のエラーは halt しない限り続行するため、エラー発生を
% ログマーカーで表明する。ランナーは TATE-PASS を数える（T8）。
\documentclass[pdflatex,paper=a5]{jlreq}
\pagestyle{empty}
\makeatletter
\begin{document}
% --- TCY-STAR-PASSTHROUGH: yoko での star 形は透過（組版結果は素の文字列）
阿伊\tatechuyoko*{AB}阿伊
\typeout{TATE-PASS: TCY-STAR-PASSTHROUGH (star form transparent in yoko)}%
% --- TCY-NONSTAR-ERROR: yoko での非 star 形は ClassError。
% エラー発生後の実行継続でマーカーを出す（エラー文言はログ検査）。
\typeout{TCY-ERROR-EXPECT}%
\tatechuyoko{12}%
\typeout{TATE-PASS: TCY-NONSTAR-ERROR-REACHED (non-star in yoko errors; run continued)}%
% --- DIR-BOXDIR-ERROR: 未知 box の方向照会は ClassError。
\newbox\unknownbox
\setbox\unknownbox=\hbox{阿}%
\typeout{BOXDIR-ERROR-EXPECT}%
\ptj@boxdir{\unknownbox}%
\typeout{TATE-PASS: DIR-BOXDIR-ERROR-REACHED (boxdir on unknown box errors; run continued)}%
\end{document}
